% SEGMENT OLP-0068-S01
% Part: first-order-logic
% Chapter: proof-systems
% Section: axiomatic-deduction

\documentclass[../../../include/open-logic-section]{subfiles}

\begin{document}

\iftag{FOL}
      {\olfileid{fol}{prf}{axd}}
      {\olfileid{pl}{prf}{axd}}

\olsection{அடிகோள் \usetoken{P}{derivation}}

அடிகோள் சார்ந்த !!{derivation}s மிகப் பழமையான, மிக எளிய தருக்க
!!{derivation} முறைமைகள். அவற்றின் !!{derivation}s வெறுமனே
!!{sentence}s இன் தொடர்கள். !!{sentence}s இன் ஒரு தொடர் சரியான
!!{derivation} ஆகக் கருதப்பட வேண்டுமெனில், அதிலுள்ள ஒவ்வொரு
!!{sentence}~$!A$ உம் பின்வரும் நிபந்தனைகளில் ஒன்றை நிறைவேற்ற வேண்டும்:
\begin{enumerate}
\item $!A$ ஓர் அடிகோள்; அல்லது
\item $!A$, கொடுக்கப்பட்ட !!{sentence}s கணம்~$\Gamma$ வின்
  !!a{element}; அல்லது
\item $!A$ ஓர் அனுமான விதியால் நியாயப்படுத்தப்படுகிறது.
\end{enumerate}
$!A$ ஓர் அடிகோளாக இருக்க, அது நிலையான !!{sentence} வார்ப்புருக்கள்
சிலவற்றில் ஒன்றின் வடிவில் இருக்க வேண்டும். முதல்-வரிசைத் தருக்கவியலுக்கான
சரித்தன்மையும் நிறைவும் கொண்ட !!{derivation} முறைமையைத் தரும் அடிகோள்
வார்ப்புருக்களின் கணங்கள் பல உள்ளன. சில கணங்கள் அவை கட்டுப்படுத்தும் தருக்க
இணைப்புகளுக்கு ஏற்ப ஒழுங்கமைக்கப்படுகின்றன. எடுத்துக்காட்டாக,
\[
!A \lif (!B \lif !A) \qquad !B \lif (!B \lor !C) \qquad (!B \land !C) \lif !B
\]
என்ற வார்ப்புருக்கள் $\lif$, $\lor$ மற்றும்~$\land$ ஐக் கட்டுப்படுத்தும்
பொதுவான அடிகோள்கள். சில அடிகோள் முறைமைகள் அடிகோள்களின் எண்ணிக்கையைக்
குறைந்தபட்சமாக்க நோக்கம்கொள்கின்றன. எந்த இணைப்புகள் முதன்மையானவையாகக்
கொள்ளப்படுகின்றன என்பதைப் பொறுத்து, ஒரே ஓர் அடிகோளைக் கொண்ட அடிகோள்
முறைமைகளையும் காண முடியும்.

% SEGMENT OLP-0068-S02
ஓர் அனுமான விதி என்பது, !!a{derivation}வில் !!a{sentence} ஒன்று
நியாயப்படுத்தப்படுவதற்கான போதிய நிபந்தனையை வழங்கும் நிபந்தனைக் கூற்று.
மோடஸ் பொனென்ஸ் அத்தகைய மிகப் பொதுவான விதி: $!A$ மற்றும் $!A \lif
!B$ ஏற்கெனவே நியாயப்படுத்தப்பட்டிருந்தால், $!B$ நியாயப்படுத்தப்படுகிறது
என்று அது கூறுகிறது. அதாவது, !!a{derivation}வில் !!{sentence}~$!B$ ஐக்
கொண்ட ஒரு வரி நியாயப்படுத்தப்பட வேண்டுமெனில், $!A$ மற்றும் $!A \lif !B$
ஆகிய இரண்டும் (ஏதோ ஒரு !!{sentence}~$!A$ க்கு)~$!B$ க்கு முன்னரே அந்த
!!{derivation}வில் இடம்பெற்றிருப்பது போதுமானது.

% SEGMENT OLP-0068-S03
அடிகோள் சார்ந்த !!{derivation}s ஐ அடிப்படையாகக் கொண்ட $\Proves$ உறவு
பின்வருமாறு வரையறுக்கப்படுகிறது: !!{sentence}~$!A$ ஐக் கடைசி வாய்பாடாகக்
கொண்ட !!a{derivation} இருந்தால், அப்போதும் அப்போது மட்டுமே,
$\Gamma \Proves !A$ (மேலுள்ள~(2) ஆல் நியாயப்படுத்தப்படும் அந்த
!!{derivation}வின் !!{sentence}s கணமாக $\Gamma$ கொள்ளப்படுகிறது).
$\Gamma$ வெறுமையாக உள்ள, அதாவது !!{derivation}வின் ஒவ்வொரு !!{sentence}வும்
(1) அல்லது~(3) ஆல் நியாயப்படுத்தப்படும், !!a{derivation} $!A$ க்கு
இருந்தால் $!A$ ஒரு தேற்றம். எடுத்துக்காட்டாக, $\Proves
!A  \lif (!B \lif (!B \lor !A))$ என்பதைக் காட்டும் !!a{derivation} இதோ:
\begin{derivation}
  1. & $!B \lif (!B \lor !A)$ \\
  2. & $(!B \lif (!B \lor !A)) \lif (!A  \lif (!B \lif (!B \lor !A)))$\\
  3. & $!A  \lif (!B \lif (!B \lor !A))$
\end{derivation}

% SEGMENT OLP-0068-S04
வரி~1 இல் உள்ள !!{sentence}, $!A \lif (!A \lor !B)$ என்ற அடிகோளின்
வடிவில் உள்ளது ($!A$ மற்றும் $!B$ ஆகியவற்றின் பங்குகள் இடமாற்றப்பட்டுள்ளன).
வரி~2 இல் உள்ள வாக்கியம் $!A \lif (!B \lif !A)$ என்ற அடிகோளின் வடிவில்
உள்ளது. ஆகவே இரு வரிகளும் நியாயப்படுத்தப்பட்டுள்ளன. வரி~3 மோடஸ்
பொனென்ஸால் நியாயப்படுத்தப்படுகிறது: அதை $!D$ என்று சுருக்கினால், வரி~2
$!C \lif !D$ என்ற வடிவில் இருக்கும்; இங்கு $!C$ என்பது $!B \lif (!B
\lor !A)$, அதாவது வரி~1.

% SEGMENT OLP-0068-S05
$\Gamma \Proves \lfalse$ எனில் கணம் $\Gamma$ முரணுடையது. நிறைவான
அடிகோள் முறைமை எந்த~$!A$ க்கும் $\lfalse \lif !A$ என்பதையும் நிறுவும்.
ஆகவே $\Gamma$ முரணுடையது எனில், எந்த~$!A$ க்கும் $\Gamma \Proves !A$.

% SEGMENT OLP-0068-S06
தருக்கவியலுக்கான அடிகோள் சார்ந்த !!{derivation} முறைமைகளை முதன்முதலில்
காட்லோப் ஃபிரேக தமது 1879 ஆம் ஆண்டின் \emph{Begriffsschrift} நூலில்
வழங்கினார்; அதனால் அந்நூல் நவீன தருக்கவியலின் முதல் படைப்பாக அடிக்கடி
கருதப்படுகிறது. ஆல்ஃபிரட் நார்த் வைட்ஹெட் மற்றும் பெர்ட்ரண்டு ரஸ்ஸல் தமது
\emph{Principia Mathematica} நூலிலும், டேவிட் ஹில்பர்ட்டும் அவரது
மாணவர்களும் 1920களிலும் அவற்றைச் செம்மைப்படுத்தினர். ஆகவே அவை அடிக்கடி
``ஃபிரேக முறைமைகள்'' அல்லது ``ஹில்பர்ட் முறைமைகள்'' எனப்படுகின்றன.
ஒரு தருக்கவியலுக்கான அடிகோள் முறைமையைக் கண்டறிவது பெரும்பாலும் எளிது;
அந்த வகையில் அவை பலவகைப் பயன்பாட்டுத் திறன் கொண்டவை. !!{derivation}s
மிக எளிய கட்டமைப்பையும் ஒன்று அல்லது இரண்டு அனுமான விதிகளையும் மட்டுமே
கொண்டிருப்பதால், அவற்றைப் \emph{பற்றிய} கருத்துகளை நிறுவுவதும்
ஒப்பீட்டளவில் எளிது. ஆனால் நடைமுறையில் அவற்றைப் பயன்படுத்துவது மிகவும்
கடினம்; அதாவது, நிறுவல்களைக் கண்டறிந்து எழுதுவது கடினம்.

\end{document}
